% GENERATED FILE. Edit canonical lessons, then run npm run generate:books.
% canonical-source-hash: fnv1a64:f200fe9b9b085cac
% canonical-lessons: ES-C384-enviar, ES-C384-subir, ES-C384-morir, ES-C384-olvidar

\chapter{Four Roads, and Where the Trail Stops}
\label{ch:enviar-y-subir}
\hlchaptermodality{384}

\begin{chapteropening}
\textbf{By the end of this chapter:} \emph{I can say enviar, subir, morir and olvidar, and name four everyday actions.}

Complete ES-C384-olvidar: name four everyday actions, and say one thing about each.
\end{chapteropening}

\section[enviar]{enviar — putting a thing on the road}

\label{lesson:ES-C384-enviar}



\begin{quote}
Say \emph{to take out}, then \emph{to look for}. (\emph{Sacar}; \emph{buscar}.)
\end{quote}

\subsection*{What to know first}
\begin{itemize}
\raggedright
  \item la calle — a road.
  \item la ciudad — where you send it.
\end{itemize}

\begin{sounds}
\begin{itemize}
\raggedright
  \item \emph{en-vi-AR}, three beats with the weight on the last. The \emph{v} is the same sound as \emph{b}, and \emph{-ia-} here is two beats rather than a glide. $\to$ pronunciation reference
\end{itemize}
\end{sounds}

\begin{cousinweb}
\textbf{enviar} — \textquotedblleft{}to send.\textquotedblright{}

Latin \textbf{via} was a road, and \textbf{inviare} was to put something \emph{on} the road. That is the whole of \emph{enviar}: sending is road-putting, and the road is still sitting in the middle of the word where you can see it.

English borrowed \emph{via} untranslated and uses it as a preposition — travelling \emph{via} Madrid is travelling by way of the road through it. \textbf{Obvious} is the same noun: \emph{ob viam}, lying in the way, so obvious things are the ones you trip over. A \textbf{viaduct} carries a road.

And an \textbf{envoy} is \emph{enviar} itself, borrowed through French: an envoy is a person who has been sent.

Now the false friend, which is a good one because the resemblance is close. \textbf{Envy} is \emph{not} in this family. It comes from Latin \emph{invidēre}, to look upon — \emph{in-} plus \emph{vidēre}, to see. To envy somebody is to look at them too hard, with the wrong sort of attention. The Romans thought of it as a kind of gaze.

So \emph{enviar} holds a road and \emph{envy} holds an eye, and the two words merely start the same way.
\end{cousinweb}

\begin{grammarlens}[title={What you've built}]
You can now name the action: \emph{enviar}. And you can find a road inside it.
\end{grammarlens}

\subsection*{Your turn}
Say these aloud:

\begin{itemize}
\raggedright
  \item \emph{enviar}
  \item \emph{enviar a la ciudad}
  \item \emph{buscar y enviar}
\end{itemize}

\begin{tcolorbox}[breakable,colback=teal!4,colframe=teal!35!black,title={Before you move on}]
Which noun is inside \emph{enviar}? (\emph{Via} — a road.) Is \emph{envy} related? (No — that is \emph{vidēre}, to see.)
\end{tcolorbox}
\section[subir]{subir — a verb for going up that starts with under}

\label{lesson:ES-C384-subir}



\begin{quote}
Say \emph{to look for}, then \emph{to send}. (\emph{Buscar}; \emph{enviar}.)
\end{quote}

\subsection*{What to know first}
\begin{itemize}
\raggedright
  \item arriba — where you end up.
  \item abajo — where you start.
\end{itemize}

\begin{sounds}
\begin{itemize}
\raggedright
  \item \emph{su-BIR}, two beats with the weight on the second. The \emph{b} between vowels is soft, made with the lips barely touching. $\to$ pronunciation reference
\end{itemize}
\end{sounds}

\begin{cousinweb}
\textbf{subir} — \textquotedblleft{}to go up.\textquotedblright{}

Look at the front of it. \textbf{Sub-} is the Latin for \emph{under}, and English uses it that way constantly: a \textbf{submarine} goes under the sea, a \textbf{subway} runs under the street, and something \textbf{subterranean} is under the ground.

So why does a verb beginning with \emph{under} mean to go \textbf{up}?

Because Latin \textbf{subire} meant to go up \emph{from underneath}. The picture is of approaching something from below and rising into it — coming up at it. The \emph{sub-} is not describing where you end up. It is describing where you started.

Once that clicks, some odd English words settle down. To \textbf{submit} is to send yourself under somebody. To \textbf{suffer} is to bear up under a weight — \emph{sub-} plus \emph{ferre}, to carry, with the \emph{b} assimilated. And \textbf{sudden} is the surprise: it comes from \emph{subitus}, which is the participle of this very verb. A sudden thing is a thing that has come up on you from below, unnoticed until it arrived.

Spanish also has \textbf{subito} as a formal word for the same idea, sitting beside the everyday \emph{subir} — the noun-shaped relative and the verb you actually use.
\end{cousinweb}

\begin{grammarlens}[title={What you've built}]
You can now name the action: \emph{subir}. And you know why it starts with \emph{under}.
\end{grammarlens}

\subsection*{Your turn}
Say these aloud:

\begin{itemize}
\raggedright
  \item \emph{subir}
  \item \emph{subir arriba}
  \item \emph{enviar y subir}
\end{itemize}

\begin{tcolorbox}[breakable,colback=teal!4,colframe=teal!35!black,title={Before you move on}]
What did \emph{sub-} mean? (Under.) So how can \emph{subir} mean to go up? (Latin \emph{subire} was to come up from underneath.)
\end{tcolorbox}
\section[morir]{morir — a pledge that dies}

\label{lesson:ES-C384-morir}



\begin{quote}
Say \emph{to send}, then \emph{to go up}. (\emph{Enviar}; \emph{subir}.)
\end{quote}

\subsection*{What to know first}
\begin{itemize}
\raggedright
  \item la noche — a word for endings.
  \item subir — the last verb you met.
\end{itemize}

\begin{sounds}
\begin{itemize}
\raggedright
  \item \emph{mo-RIR}, two beats with the weight on the second. Both \emph{r} sounds are single taps here. $\to$ pronunciation reference
\end{itemize}
\end{sounds}

\begin{cousinweb}
\textbf{morir} — \textquotedblleft{}to die.\textquotedblright{}

Latin \textbf{mori}, to die, with \textbf{mors} the noun. English took the family whole: \textbf{mortal}, \textbf{mortality}, \textbf{mortuary}, \textbf{immortal}.

The one worth stopping on is \textbf{mortgage}, because almost nobody reads it. It is two French words: \emph{mort} + \emph{gage}, a \textbf{dead pledge}. The medieval lawyers meant something precise by it. The pledge is alive while the debt runs; when the debt is paid the pledge dies, and the property comes free. So a mortgage is named for its own ending.

Now a careful point, because there is a tempting shortcut here that is wrong. \textbf{Murder} does belong to the same family — but not through Latin. It is a native English word, Old English \emph{morthor} — written with a letter English has since lost — arriving through Germanic. What it shares with \emph{mori} is the far older Indo-European root behind both. Calling \emph{murder} a Latin word would be false; calling it a cousin that travelled a different road is right, and the distinction is the kind that separates a real etymology from a plausible one.

One more to refuse: \textbf{morgue} is not in this family. It came into French with a different meaning entirely and was only later attached to a place for the dead.
\end{cousinweb}

\begin{grammarlens}[title={What you've built}]
You can now name the action: \emph{morir}. And you can read \emph{mortgage} as two words.
\end{grammarlens}

\subsection*{Your turn}
Say these aloud:

\begin{itemize}
\raggedright
  \item \emph{morir}
  \item \emph{subir y morir}
  \item \emph{enviar}
\end{itemize}

\begin{tcolorbox}[breakable,colback=teal!4,colframe=teal!35!black,title={Before you move on}]
What are the two halves of \emph{mortgage}? (\emph{Mort} — dead; \emph{gage} — pledge.) Did \emph{murder} come through Latin? (No — through Germanic, from the same older root.)
\end{tcolorbox}
\section[olvidar]{olvidar — where to stop}

\label{lesson:ES-C384-olvidar}



\begin{quote}
Say \emph{to go up}, then \emph{to die}. (\emph{Subir}; \emph{morir}.)
\end{quote}

\subsection*{What to know first}
\begin{itemize}
\raggedright
  \item buscar — what you do after forgetting.
  \item la llave — a thing worth not forgetting.
\end{itemize}

\begin{sounds}
\begin{itemize}
\raggedright
  \item \emph{ol-vi-DAR}, three beats with the weight on the last. The \emph{v} is the same sound as \emph{b}, and the \emph{d} between vowels is soft, close to the \emph{th} of \emph{this}. $\to$ pronunciation reference
\end{itemize}
\end{sounds}

\begin{cousinweb}
\textbf{olvidar} — \textquotedblleft{}to forget.\textquotedblright{}

Latin \textbf{oblivisci} meant to forget. Spoken Latin rebuilt it as \emph{oblitare}, a tidier verb with regular endings, and that is what came down into Spanish as \emph{olvidar}.

English has the noun: \textbf{oblivion}, the state of being forgotten, and \textbf{oblivious}, unaware. Same word, different shape.

And here the lesson does something the others have not, which is stop.

There is a further step everybody wants to take. \emph{Oblivisci} looks like it comes apart into \emph{ob-} and a root meaning smooth — as though forgetting were smoothing something away, wiping a wax tablet flat. It is a lovely picture. It is also one that specialists in Latin's own history have examined and rejected, on the evidence of how the word behaves rather than on how good the story sounds.

So the trail runs: \emph{oblivisci}, \emph{oblitare}, \emph{olvidar}, and beside it English \emph{oblivion}. Past that point, nobody can take you.

That is worth practising as much as any of the roads that go all the way through. A word's history is a piece of evidence, and evidence runs out. Knowing where it runs out is part of knowing the word.
\end{cousinweb}

\begin{grammarlens}[title={What you've built}]
You can now name four more everyday actions: \emph{enviar}, \emph{subir}, \emph{morir}, \emph{olvidar}. Say one thing about each.
\end{grammarlens}

\subsection*{Your turn}
Say these aloud:

\begin{itemize}
\raggedright
  \item \emph{olvidar}
  \item \emph{olvidar la llave}
  \item \emph{enviar, subir, morir, olvidar}
\end{itemize}

\begin{tcolorbox}[breakable,colback=teal!4,colframe=teal!35!black,title={Before you move on}]
Which English noun is the same word? (\emph{Oblivion}.) And why does the lesson stop there? (The next step is rejected by specialists.)
\end{tcolorbox}
